%%
% This is an Overleaf template for presentations
% using the TUM Corporate Desing https://www.tum.de/cd
%
% For further details on how to use the template, take a look at our
% GitLab repository and browse through our test documents
% https://gitlab.lrz.de/latex4ei/tum-templates.
%
% The tumbeamer class is based on the beamer class.
% If you need further customization please consult the beamer class guide
% https://ctan.org/pkg/beamer.
% Additional class options are passed down to the base class.
%
% If you encounter any bugs or undesired behaviour, please raise an issue
% in our GitLab repository
% https://gitlab.lrz.de/latex4ei/tum-templates/issues
% and provide a description and minimal working example of your problem.
%%

\PassOptionsToClass{onlytextwidth}{beamer}

\documentclass[
  german,            % define the document language (english, german)
  aspectratio=169,    % define the aspect ratio (169, 43)
  % handout=2on1,       % create handout with multiple slides (2on1, 4on1)
  % partpage=false,     % insert page at beginning of parts (true, false)
  % sectionpage=true,   % insert page at beginning of sections (true, false)
]{tumbeamer}


% load additional packages
\usepackage{tikz}
\usepackage{circuitikz}
\usepackage{url}
\usepackage{hyperref}
\usepackage{pgf}
\usepackage{pgfplots}
\usepackage{babel}[ngerman]
\usepackage{csquotes}[autostyle]
\usepackage[useregional]{datetime2}
\usepackage{float}
\usepackage{graphicx}
\usepackage{amsmath}
\usepackage{amssymb}
\usepackage{xcolor}
\usepackage[cache=true]{minted}
\usemintedstyle{borland}
%\usepackage{listings}
%\usepackage{tikz-timing}
%\usepackage{ifthen}
\usepackage{ulem}
\usepackage{enumerate}
\usepackage{karnaugh-map}
%\usepackage{forest}
\usepackage{pifont}


% tikz
\usetikzlibrary{overlay-beamer-styles}
\usetikzlibrary{arrows,backgrounds,positioning,shapes,patterns,patterns.meta,matrix,arrows.meta,shapes.geometric}
\usetikzlibrary{matrix, fit, calc}
%\usetikzlibrary{automata}

% requires circuitikz >= 1.1.0
% for distros with older distributions, install TeX Live manually
% instead of using your package manager
% see: https://tug.org/texlive/quickinstall.html
\ctikzset{logic ports=ieee}

% minted
\setminted{
    fontsize=\small, 
    frame=none,
    breaklines=false,
}

% lstlisting
% \lstset{escapeinside={<@}{@>}}

% image path
\graphicspath{ {./resources/} }

% commands
\newcommand{\n}[1]{\overline{#1}}

% https://tex.stackexchange.com/a/7045
\newcommand*\circled[1]{\tikz[baseline=(char.base)]{
		\node[shape=circle,draw,inner sep=2pt, font=\scriptsize] (char) {#1};}}
	
\newcommand*\colorcirc[1]{\tikz[baseline=(char.base)]{
			\node[shape=circle,fill=#1,inner sep=2pt] {};}}
		
\newcommand{\cmark}{\ding{51}}%
\newcommand{\xmark}{\ding{55}}%

% beamer
\setbeamercolor{footnote}{fg=black}
\setbeamercolor{footnote mark}{fg=black}

% presentation metadata
\title{Übung 14: Fragestunde}
\subtitle{Einführung in die Rechnerarchitektur}
\author{Niklas Ladurner}

\institute{\theChairName\\\theDepartmentName\\\theUniversityName}
\date{\DTMdisplaydate{2025}{01}{31}{-1}}

\footline{\insertauthor~|~\insertshorttitle~|~\insertshortdate}


% macro to configure the style of the presentation
\TUMbeamersetup{
  title page = TUM tower,         % style of the title page
  part page = TUM toc,            % style of part pages
  section page = TUM toc,         % style of section pages
  content page = TUM more space,  % style of normal content pages
  tower scale = 1.0,              % scaling factor of TUM tower (if used)
  headline = TUM threeliner,      % which variation of headline to use
  footline = TUM default,         % which variation of footline to use
  % configure on which pages headlines and footlines should be printed
  headline on = {title page},
  footline on = {every page, title page=false},
}

\setbeamercolor{enumerate item}{fg=black}

\begin{document}

\maketitle

\begin{frame}[c]{Feedback}{}
	\begin{columns}[c]
		\begin{column}{0.5\textwidth}
			\begin{center}
				\LARGE  \href{https://t1p.de/era2425}{t1p.de/era2425}\\
				\includegraphics[height=0.5\textheight]{feedback_qr.png}
			\end{center}
		\end{column}
		\begin{column}{0.5\textwidth}
			\begin{center}
				\LARGE  \href{https://home.in.tum.de/~ladu/}{home.in.tum.de/\string~ladu/}\\
				\includegraphics[height=0.5\textheight]{website_qr.png}
			\end{center}
		\end{column}
	\end{columns}
\end{frame}

\begin{frame}[c, fragile]{}{}
	\begin{center}
		\LARGE  Keine Garantie für die Richtigkeit der Tutorfolien.

		\Large Bei Unklarheiten/Unstimmigkeiten haben VL/ZÜ-Folien recht!
	\end{center}
\end{frame}

\begin{frame}[c]{Unterlagen aus den Tutorübungen}{}
	\begin{columns}[c]
		\begin{column}{0.6\textwidth}
			\begin{itemize}
				\item Slides/Mitschriften gerne mit Kommilitonen teilen
				\item Bei Weitergabe von Dokumenten habe ich keine Möglichkeit, Fehler zu korrigieren
				\item Besser: Link zur Website teilen
				\item \underline{KEIN UPLOAD} auf Studocu/Studydrive oder sonst irgendwo
				\item Slides von Bjarne Hansen: \href{https://home.cit.tum.de/~hbj}{home.in.tum.de/\string~hbj/}
			\end{itemize}
		\end{column}
		\begin{column}{0.4\textwidth}
			\begin{center}
				\scalebox{-1}[1]{\includegraphics[width=0.5\textwidth]{w14_popcat.png}}
			\end{center}
		\end{column}
	\end{columns}

\end{frame}

\begin{frame}[c]{Klausur: Dos and Don'ts}{}
	\begin{columns}[c]
		\begin{column}{0.5\textwidth}
			\centering \textbf{Dos} \textcolor{green}{\cmark}
			\vspace{\baselineskip}
			\begin{itemize}
				\item Student Card unbedingt mitbringen
				\item genügend Kugelschreiber mitnehmen
				\item Essen/Trinken mitnehmen
				\item ca. 20 Minuten vorher im Hörsaal sein
				\item Wörterbuch Deutsch $\leftrightarrow$ Muttersprache mitnehmen, falls benötigt
			\end{itemize}
		\end{column}
		\begin{column}{0.5\textwidth}
			\centering \textbf{Don'ts} \textcolor{red}{\xmark}
			\vspace{\baselineskip}
			\begin{itemize}
				\item rote oder grüne Stifte, Bleistift
				\item Tipp-Ex oder wärmeempfindliche Stifte
				\item Taschenrechner (zählt als Unterschleif)
				\item eigene Schmierblätter
				\item irgendwelche Spickzettel -- wir finden die sowieso
			\end{itemize}
		\end{column}
	\end{columns}
\end{frame}

\begin{frame}[c]{Klausur: Vorbereitung}{}
	\begin{itemize}
		\item Übungsblätter rechnen
		\item Probeklausur rechnen, um Gefühl für Klausurformat zu bekommen. Am besten ohne Hilfsmittel/Lösung + Stoppuhr
		\item Assembly üben! Aufgaben oft sehr einfach
		\item zusätzliche Assembly-Aufgaben: \href{https://github.com/ladnik/RISC-V-Assembly-Introductory-Exercises}{github.com/ladnik}
		\item meist ist ERA eine Überhangsklausur: zuerst Aufgaben machen, die ihr gut könnt
	\end{itemize}
\end{frame}


\begin{frame}[fragile, c]{Statistik zur Endterm WS 23/24}{}
	\begin{columns}[c]
		\begin{column}{\textwidth}
			\begin{tikzpicture}
				\definecolor{col10}{RGB}{0, 255, 0}
				\definecolor{col13}{RGB}{51, 255, 0}
				\definecolor{col17}{RGB}{102, 255, 0}
				\definecolor{col20}{RGB}{154, 255, 0}
				\definecolor{col23}{RGB}{205, 255, 0}
				\definecolor{col27}{RGB}{255, 255, 0}
				\definecolor{col30}{RGB}{255, 205, 0}
				\definecolor{col33}{RGB}{255, 154, 0}
				\definecolor{col37}{RGB}{255, 102, 0}
				\definecolor{col40}{RGB}{255, 51, 0}
				\definecolor{col43}{RGB}{255, 0, 0}
				\definecolor{col47}{RGB}{170, 0, 0}
				\definecolor{col50}{RGB}{85, 0, 0}

				\begin{axis}[
						scale only axis,
						axis line style={draw=none},
						tick style={draw=none},
						width=\textwidth,
						height=0.75\textheight,
						yticklabel=\empty,
						xtick={1.0, 1.3, 1.7, 2.0, 2.3, 2.7, 3.0, 3.3, 3.7, 4.0, 4.3, 4.7, 5.0},
						x tick label style={
							/pgf/number format/.cd,
							fixed,
							fixed zerofill,
							precision=1,
							/tikz/.cd
						},
						nodes near coords,
						bar width=0.8cm,
					]
					\addplot[ybar, fill=col10] coordinates {(1.0, 37)};
					\addplot[ybar, fill=col13] coordinates {(1.3, 34)};
					\addplot[ybar, fill=col17] coordinates {(1.7, 37)};
					\addplot[ybar, fill=col20] coordinates {(2.0, 42)};
					\addplot[ybar, fill=col23] coordinates {(2.3, 59)};
					\addplot[ybar, fill=col27] coordinates {(2.7, 59)};
					\addplot[ybar, fill=col30] coordinates {(3.0, 62)};
					\addplot[ybar, fill=col33] coordinates {(3.3, 46)};
					\addplot[ybar, fill=col37] coordinates {(3.7, 73)};
					\addplot[ybar, fill=col40] coordinates {(4.0, 33)};
					\addplot[ybar, fill=col43] coordinates {(4.3, 51)};
					\addplot[ybar, fill=col47] coordinates {(4.7, 90)};
					\addplot[ybar, fill=col50] coordinates {(5.0, 195)};
					
					\node [align=left, anchor=west] at (axis cs:0.8,150) {
						\fontfamily{cmr}\selectfont Angetreten:  818\\
						\fontfamily{cmr}\selectfont Anteil negativer Beurteilungen: 41\% \\
						\fontfamily{cmr}\selectfont Durchschnitt gesamt: 3.52\\
						\fontfamily{cmr}\selectfont Durchschnitt bestanden:  2.62\\
					};    
				\end{axis}
			\end{tikzpicture}
		\end{column}
	\end{columns}
\end{frame}

\begin{frame}[fragile, c]{Statistik zur Retake WS 23/24}{}
	\begin{columns}[c]
		\begin{column}{\textwidth}
			\begin{tikzpicture}
				\definecolor{col10}{RGB}{0, 255, 0}
				\definecolor{col13}{RGB}{51, 255, 0}
				\definecolor{col17}{RGB}{102, 255, 0}
				\definecolor{col20}{RGB}{154, 255, 0}
				\definecolor{col23}{RGB}{205, 255, 0}
				\definecolor{col27}{RGB}{255, 255, 0}
				\definecolor{col30}{RGB}{255, 205, 0}
				\definecolor{col33}{RGB}{255, 154, 0}
				\definecolor{col37}{RGB}{255, 102, 0}
				\definecolor{col40}{RGB}{255, 51, 0}
				\definecolor{col43}{RGB}{255, 0, 0}
				\definecolor{col47}{RGB}{170, 0, 0}
				\definecolor{col50}{RGB}{85, 0, 0}
				
				\begin{axis}[
					scale only axis,
					axis line style={draw=none},
					tick style={draw=none},
					ymin=0,
					width=\textwidth,
					height=0.75\textheight,
					yticklabel=\empty,
					xtick={1.0, 1.3, 1.7, 2.0, 2.3, 2.7, 3.0, 3.3, 3.7, 4.0, 4.3, 4.7, 5.0},
					x tick label style={
						/pgf/number format/.cd,
						fixed,
						fixed zerofill,
						precision=1,
						/tikz/.cd
					},
					nodes near coords,
					bar width=0.8cm,
					]
					\addplot[ybar, fill=col10] coordinates {(1.0, 2)};
					\addplot[ybar, fill=col13] coordinates {(1.3, 4)};
					\addplot[ybar, fill=col17] coordinates {(1.7, 9)};
					\addplot[ybar, fill=col20] coordinates {(2.0, 6)};
					\addplot[ybar, fill=col23] coordinates {(2.3, 7)};
					\addplot[ybar, fill=col27] coordinates {(2.7, 16)};
					\addplot[ybar, fill=col30] coordinates {(3.0, 23)};
					\addplot[ybar, fill=col33] coordinates {(3.3, 25)};
					\addplot[ybar, fill=col37] coordinates {(3.7, 33)};
					\addplot[ybar, fill=col40] coordinates {(4.0, 44)};
					\addplot[ybar, fill=col43] coordinates {(4.3, 32)};
					\addplot[ybar, fill=col47] coordinates {(4.7, 20)};
					\addplot[ybar, fill=col50] coordinates {(5.0, 72)};
					
					\node [align=left, anchor=west] at (axis cs:0.8,54.25) {
						\fontfamily{cmr}\selectfont Angetreten:  293\\
						\fontfamily{cmr}\selectfont Anteil negativer Beurteilungen: 42\% \\
						\fontfamily{cmr}\selectfont Durchschnitt gesamt: 3.87\\
						\fontfamily{cmr}\selectfont Durchschnitt bestanden:  3.22\\
					};    
				\end{axis}
			\end{tikzpicture}
		\end{column}
	\end{columns}
\end{frame}

\begin{frame}[c]{}{}
	\begin{center}
		\LARGE Fragen?
	\end{center}
\end{frame}


\begin{frame}[c]{}{}
	\begin{center}
		\LARGE Viel Erfolg bei der Klausur :)
	\end{center}
\end{frame}

\begin{frame}[c, fragile]{Links}{}
	\begin{itemize}
		\item Zulip: \href{https://zulip.in.tum.de/#narrow/stream/2661-ERA-Tutorium---Do-1600-1}{\enquote{ERA Tutorium - Do-1600-1}}
		      bzw. \href{https://zulip.in.tum.de/#narrow/stream/2675-ERA-Tutorium---Fr-1500-2 }{\enquote{ERA Tutorium - Fr-1500-2}}
		\item \href{https://www.moodle.tum.de/course/view.php?id=100633}{ERA-Moodle-Kurs}
		\item \href{https://artemis.in.tum.de/courses/401}{ERA-Artemis-Kurs}
	\end{itemize}
\end{frame}

\maketitle

\end{document}
